\documentclass[14pt,a4paper]{extarticle}
\usepackage[T2A]{fontenc}\usepackage[utf8]{inputenc}\usepackage[russian]{babel}
\usepackage{geometry}\geometry{left=30mm,right=15mm,top=20mm,bottom=20mm}
\setlength{\parindent}{1.25cm}\setlength{\parskip}{0.3em}\sloppy
\begin{document}\begin{titlepage}\begin{center}
САНКТ-ПЕТЕРБУРГСКИЙ НАЦИОНАЛЬНЫЙ\\ ИССЛЕДОВАТЕЛЬСКИЙ УНИВЕРСИТЕТ ИТМО
\vspace{1.5cm}

Дисциплина: Бэк-энд разработка
\vspace{1.5cm}

\textbf{Отчёт}\\[0.5cm]
Домашняя работа 10\\[0.5cm]
\textbf{Сценарии нагрузочного тестирования}
\vfill
\begin{flushright}Выполнил:\\ Цатинян Артём\\ Группа БР1.2\\[0.8cm] Проверил:\\ Добряков Д. И.\end{flushright}
\vfill Санкт-Петербург\\ 2026 г.\end{center}\end{titlepage}
\section{Задача и состояние}
Спроектировать нагрузочные сценарии, метрики и целевые показатели для приложения. Численные значения ниже являются учебными целями, а не результатами испытаний или согласованными обязательствами.
\section{Сценарии}
Для инструмента k6 предлагаются сценарии: чтение каталога с фильтрами, авторизация тестового пользователя, создание и просмотр бронирования, получение отзывов. Запросы должны использовать изолированные тестовые данные. Точные пути и тела запросов проверяются по OpenAPI перед выполнением.
\section{Профили нагрузки}
Базовый профиль последовательно увеличивает интенсивность, удерживает нагрузку и снижает её; стресс-профиль повышает интенсивность до появления признаков насыщения; длительный профиль проверяет устойчивость при стабильной нагрузке. Параметры подбираются по ресурсам стенда и согласуются до запуска.
\section{Метрики и цели}
Основные метрики: количество запросов в секунду, доля ошибок, p95 и p99 времени ответа, загрузка CPU, памяти и БД. Пример предварительного ориентира для операций чтения: p95 менее 500 мс и доля ошибок менее 1\%; до измерений эти значения не считаются достигнутыми. Для изменения данных и авторизации цели следует согласовать отдельно.
\section{Методика}
Перед запуском подготовить данные и мониторинг, провести прогрев, зафиксировать конфигурацию и интервалы, затем сохранить агрегированные результаты и ошибки. Повторить тест после изменений при одинаковом профиле. Не направлять нагрузку на рабочую среду без согласования.
\section{Вывод}
Составлен план сценариев и измерений; скрипты и фактическое тестирование остаются невыполненными.
\end{document}
